\section{Technical Promise Versus Standard Gravity}

New languages emerge when an existing compression package stops matching the dominant pain. They fracture around runtime ceilings, ecosystem gaps, tooling weakness, migration cost, hiring constraints, interop friction, and fashion that lacks enforcement. A language can be technically excellent and still fail as a default company architecture if it lacks boring deployment, security scanning, observability, package maturity, database migration norms, and enough public examples for agents and humans to learn from.

Julia is the cleanest example of legitimate promise without universal default status. Its multiple dispatch, JIT compilation, and scientific-computing focus attack the two-language problem in numerical work: prototype in a high-abstraction language, then rewrite hot paths in C, C++, or Fortran \cite{bezansonJulia2017}. That is a real problem, and Julia remains a serious answer for many scientific domains. The adoption lesson is not that Julia is bad. The lesson is that solving an inner-loop expression problem does not automatically solve company-scale proof routing, repair locality, cloud operations, security review, hiring, packaging, and cross-language contract discipline.

\begin{table}[t]
\caption{Technical promise versus standard gravity.}
\label{tab:standard-gravity}
\centering
\scriptsize
\begin{tabularx}{\columnwidth}{L{0.24\columnwidth} Y}
\toprule
\JKTableHeader
\textbf{Ecosystem} & \textbf{Jankurai interpretation} \\
\midrule
Julia & Real scientific-computing promise; broader product proof routing and deployment defaults remain profile-specific \cite{bezansonJulia2017}. \\
Elm/F\#/Clojure & Strong reliability or functional-design lessons; narrower hiring/corpus surfaces require explicit conventions. \\
Haskell/Scala & Proof power helps only when abstraction and build conventions are common enough for repair agents. \\
Elixir/Phoenix & Excellent realtime profile when equivalent proof receipts, owner routes, and witnesses exist. \\
Nim/Crystal/D & Ergonomics and native compilation are not enough without ecosystem-depth and CI/security defaults. \\
\bottomrule
\end{tabularx}
\end{table}

The shelf is not a graveyard of bad ideas. It is evidence that technical elegance and default-standard status are different things. Jankurai Core does not choose a language. It asks whether the chosen language and ecosystem can emit owner routes, proof receipts, generated-zone evidence, and merge witnesses cheaply enough to reduce drift over time.

This distinction matters because stack preference can become another review-subjectivity gap. A Rust, Go, JVM, .NET, Elixir, Rails, or TypeScript-heavy repository can all be easy or hard for agents depending on whether the repo exposes deterministic setup, owner cells, generated boundaries, stable errors, and proof receipts. The reference profile is therefore adoption guidance, not a stack war. Agent-first design is the portable requirement: make the repository shape explain itself before the reviewer or model has to guess.
